\documentclass[a4paper]{article}

\usepackage[fontset=fandol]{ctex}
\usepackage{amsmath, amssymb}
\usepackage{graphicx}
\usepackage[margin=2.5cm]{geometry}
\usepackage[most]{tcolorbox}
\usepackage{etoolbox}
\usepackage{listings}
\usepackage{hyperref}
\usepackage{booktabs}
\usepackage{subcaption}
\usepackage{float}
\usepackage{tikz}
\IfFileExists{pgfplots.sty}{
    \usepackage{pgfplots}
    \pgfplotsset{compat=1.18}
}{}

\newtcolorbox{knowledgebox}[1]{
    enhanced, colback=blue!5!white, colframe=blue!75!black, colbacktitle=blue!75!black,
    coltitle=white, fonttitle=\bfseries, title=#1, attach boxed title to top left={yshift=-2mm, xshift=2mm},
    boxrule=1pt, sharp corners
}

\newtcolorbox{importantbox}[1]{
    enhanced, colback=yellow!10!white, colframe=yellow!80!black, colbacktitle=yellow!80!black,
    coltitle=black, fonttitle=\bfseries, title=#1, sharp corners
}

\newtcolorbox{warningbox}[1]{
    enhanced, colback=red!5!white, colframe=red!75!black, colbacktitle=red!75!black,
    coltitle=white, fonttitle=\bfseries, title=#1, sharp corners
}

\lstset{
    language=Python,
    basicstyle=\ttfamily\small,
    keywordstyle=\color{blue},
    stringstyle=\color{red!60!black},
    commentstyle=\color{green!60!black},
    breaklines=true,
    frame=single,
    numbers=left,
    numberstyle=\tiny\color{gray},
    captionpos=b,
    extendedchars=false
}

\newcommand{\notetitle}{CME 295 课程笔记：Lecture 3 Large Language Models}
\newcommand{\noteauthors}{Codex}
\newcommand{\notedate}{2026-04-13}
\newcommand{\videochannel}{Stanford Online}
\newcommand{\videopublishdate}{2025-10-17}
\newcommand{\videoduration}{1:48:45}
\newcommand{\videourl}{https://www.youtube.com/watch?v=Q5baLehv5So}
\newcommand{\videocoverpath}{cover.jpg}

\begin{document}

\begin{titlepage}
\centering
{\Large 课程笔记\par}
\vspace{1.2cm}
{\huge\bfseries \notetitle\par}
\vspace{0.8cm}
{\large \noteauthors\par}
\vspace{0.3cm}
{\large \notedate\par}
\vspace{1.2cm}

\ifdefempty{\videocoverpath}{
{\small 请在生成笔记时填入视频封面图的本地路径。\par}
}{
\includegraphics[width=0.82\textwidth,height=0.45\textheight,keepaspectratio]{\videocoverpath}\par
}

\vfill
\begin{tcolorbox}[width=0.9\textwidth, colback=black!2!white, colframe=black!60, sharp corners]
\textbf{视频作者/频道}：\videochannel\par
\textbf{发布日期}：\videopublishdate\par
\textbf{视频时长}：\videoduration\par
\textbf{视频链接}：
\href{\videourl}{\nolinkurl{\videourl}}
\end{tcolorbox}
\end{titlepage}

\tableofcontents
\newpage

\section{什么是 LLM：定义、尺度与主干结构}

第三讲正式进入大语言模型本体。课程先用上一讲的三类 Transformer-based 模型做 recap，然后把范围收缩到今天最主流的一类：decoder-only、自回归、text-to-text 的大规模语言模型。

讲者对 ``large'' 的定义非常工程化，不是模糊形容词，而是至少包含三层含义：
\begin{itemize}
\item 参数规模大，通常从十亿级往上走；
\item 训练数据量大，常见为数百亿到万亿 token；
\item 训练和推理所需算力大，需要大量 GPU 资源和复杂系统优化。
\end{itemize}

\begin{figure}[H]
\centering
\includegraphics[width=0.72\textwidth]{pdf_pages/page-13.png}
\caption{课程强调现代 LLM 的主干几乎都是 decoder-only Transformer。}
\end{figure}

课程还特别澄清了一个常见误解：不是所有“大模型”都算今天语境里的 LLM。按照本讲的定义，BERT 这类 encoder-only 模型虽然参数不小，但不做自回归文本生成，因此通常不被归入现代 LLM 的主流范畴。

\begin{knowledgebox}{本讲默认的 LLM 定义}
一个现代 LLM 通常是 decoder-only Transformer，通过 next-token prediction 训练，在超大语料和超大算力下学到能够做文本生成的条件分布模型。
\end{knowledgebox}

\subsection{本章小结}

LLM 首先是 language model，其次才是 large。课程之所以从结构、数据和算力三个维度同时定义它，是为了说明后面所有技巧都围绕这三个量级约束展开。

\section{MoE：为什么不必每次都激活全部参数}

\subsection{从“所有参数一起算”到“只选部分专家”}

如果一个模型有海量参数，那么每次前向都让所有参数参与计算，会带来极高成本。课程用“房间里有数学家、物理学家、化学家和历史学家，你问一个数学问题是否需要所有人一起回答”这个比喻引出 Mixture of Experts（MoE）。

\begin{figure}[H]
\centering
\includegraphics[width=0.78\textwidth]{pdf_pages/page-23.png}
\caption{MoE 的核心结构：router/gate 根据输入选择少量专家参与计算。}
\end{figure}

在 MoE 中，给定输入 \(x\)，会有一个 gate 或 router 决定哪些 expert 需要被激活。于是模型的“总参数量”可以很大，但“单次激活参数量”相对较小，从而提高参数规模与推理成本之间的比例。

\subsection{Dense MoE 与 Sparse MoE}

课程区分了两类 MoE：
\begin{itemize}
\item Dense MoE：所有 expert 都参与，只是权重不同；
\item Sparse MoE：只选极少数 expert，真正做到稀疏激活。
\end{itemize}

现代大模型讨论的 MoE，主要指 sparse MoE。它的价值就在于：总参数可以继续扩张，但每个 token 实际只路由到少数专家，不必让全部参数同时上线。

\subsection{MoE 放在 Transformer 的哪里}

课程说明，MoE 通常不是替换 attention，而是替换或增强 FFN 子层。理由很直观：attention 负责 token 间的信息交互，而 FFN 更像逐位置的特征变换模块，把 expert 机制放在这里更自然。

\subsection{路由塌缩与负载均衡}

MoE 的难点不在概念，而在训练。如果 router 总是把大多数 token 分给同一个 expert，就会出现 routing collapse，导致专家利用率极不平衡、学习效率很差。

\begin{figure}[H]
\centering
\includegraphics[width=0.70\textwidth]{pdf_pages/page-29.png}
\caption{课程把 routing collapse 定义为大多数 token 总是流向同一专家。}
\end{figure}

因此训练时通常要加入负载均衡目标，让不同 expert 获得更均匀的 token 流量。课程在口头说明里也强调，MoE 路由通常是在 token 级别发生的，即同一序列中的不同 token 完全可能流向不同专家。

\begin{figure}[H]
\centering
\includegraphics[width=0.68\textwidth]{pdf_pages/page-31.png}
\caption{按 token 着色的专家路由图说明，不同 token 会被送往不同 expert。}
\end{figure}

\subsection{本章小结}

MoE 的核心不是“多几个模块”，而是把模型扩张方式从“所有参数永远一起工作”改成“有条件地激活部分参数”。它让参数规模继续上升成为可能，但代价是必须解决路由训练、负载均衡和系统实现问题。

\section{从 next-token prediction 到具体解码策略}

\subsection{语言模型真正学到的是什么}

课程从最基本的 next-token prediction 出发，强调语言模型学习的是条件分布：
\[
p_\theta(x_t \mid x_{<t})
\]

其中：
\begin{itemize}
\item \(x_{<t}\)：当前位置之前的 token 前缀；
\item \(x_t\)：当前位置候选 token；
\item \(p_\theta\)：模型参数为 \(\theta\) 的条件概率分布。
\end{itemize}

训练时模型学的是这个分布；推理时真正困难的，是如何从这个分布中选 token。课程用 greedy、beam search 与 sampling 三种策略逐步展开。

\begin{figure}[H]
\centering
\includegraphics[width=0.78\textwidth]{pdf_pages/page-45.png}
\caption{Greedy decoding 每一步都选当前概率最高的 token。}
\end{figure}

\subsection{Greedy decoding 与 Beam search}

Greedy decoding 最简单，但只追求局部最优，容易在整体序列层面不自然或不够多样。Beam search 则保留多个候选路径，尝试避免局部最优陷阱。

\begin{figure}[H]
\centering
\includegraphics[width=0.78\textwidth]{pdf_pages/page-47.png}
\caption{Beam search 通过维护多个候选前缀，寻找更高整体概率的序列。}
\end{figure}

课程也提醒 beam search 并不是今天开放式生成的默认选择，因为它更适合机器翻译等追求高条件概率序列的任务，而开放式对话、创作往往更需要多样性和自然性。

\subsection{Sampling、Top-k 与 Top-p}

为了保留多样性，生成系统常用 sampling：不是永远取最高概率 token，而是按分布采样。进一步地：
\begin{itemize}
\item Top-k：只保留前 \(k\) 个最可能 token 后再采样；
\item Top-p：保留累计概率达到 \(p\) 的最小 token 集合后再采样。
\end{itemize}

\begin{figure}[H]
\centering
\includegraphics[width=0.80\textwidth]{pdf_pages/page-56.png}
\caption{Sampling 不再固定选择最大概率 token，而是在概率分布上抽样。}
\end{figure}

这两种方法都在做同一件事：把极低概率、往往会导致胡言乱语的长尾 token 剪掉，同时保留一定随机性。

\subsection{Temperature 与 guided decoding}

temperature 用来调节分布尖锐程度：
\[
p_i^{(T)} = \frac{\exp(z_i/T)}{\sum_j \exp(z_j/T)}
\]

符号说明如下：
\begin{itemize}
\item \(z_i\)：第 \(i\) 个 token 的 logit；
\item \(T\)：温度；
\item \(T<1\) 时分布更尖锐，更保守；
\item \(T>1\) 时分布更平坦，更发散。
\end{itemize}

\begin{figure}[H]
\centering
\includegraphics[width=0.74\textwidth]{pdf_pages/page-63.png}
\caption{温度并不改变 token 排序本身，而是改变概率分布的尖锐程度。}
\end{figure}

guided decoding 则是另一条路线：不是调分布，而是硬性约束哪些 token 合法。课程用生成 JSON 的例子说明，如果输出必须满足某种格式，那么可以在每一步只允许“语法上合法”的 token 进入候选集合。

\begin{figure}[H]
\centering
\includegraphics[width=0.78\textwidth]{pdf_pages/page-65.png}
\caption{Guided decoding 通过约束合法 token 集，强制输出满足预定结构。}
\end{figure}

\subsection{本章小结}

模型训练学的是条件分布，推理系统解决的是“如何从分布中取样”。Greedy、beam、sampling、temperature 和 guided decoding 都在处理这一层。生成质量不仅取决于模型本身，也取决于解码策略。

\section{上下文窗口与提示策略：从 ICL 到 CoT}

\subsection{Context length 不只是一个数字}

课程把 context length 定义为模型在单次前向中能看到的输入长度。它决定了提示词、检索文档、历史对话和中间推理痕迹能否同时放进上下文窗口。随着应用复杂化，context length 直接限制了模型能做多长链条的条件推理。

\subsection{In-Context Learning}

ICL（In-Context Learning）指的是不改模型参数，只通过 prompt 中给出的指令和示例，让模型在当前上下文里“表现得像学会了任务”。

\begin{figure}[H]
\centering
\includegraphics[width=0.80\textwidth]{pdf_pages/page-81.png}
\caption{ICL 通过在 prompt 里放入任务说明与示例，让模型在上下文中完成适配。}
\end{figure}

课程把 zero-shot、one-shot、few-shot 都放在这一框架里理解：区别只在 prompt 里是否给例子、给几个例子，而不是模型内部结构发生变化。

\subsection{Chain of Thought 与 Self-Consistency}

Chain of Thought（CoT）的核心观察是：对于复杂推理任务，让模型显式生成中间推理步骤，往往比直接输出最终答案更稳。

\begin{figure}[H]
\centering
\includegraphics[width=0.78\textwidth]{pdf_pages/page-84.png}
\caption{CoT 鼓励模型先写出中间推理过程，再给出最终答案。}
\end{figure}

但 CoT 也会带来更长输出和更高成本，所以课程进一步引出 self-consistency：与其只走一条推理链，不如采样多条 reasoning path，再对最终答案做聚合。这样做的逻辑是，单条链可能偶然走偏，但多个独立样本的共同结论往往更稳。

\subsection{本章小结}

Prompting 技术的共同点是：不改参数、改输入。ICL 解决快速适配，CoT 解决复杂推理可分解性，self-consistency 则用多路径聚合降低单次采样的偶然性。这些策略解释了为什么“会写 prompt”在 LLM 时代成为真实能力。

\section{推理优化：KV cache、PagedAttention 与更快解码}

\subsection{为什么推理会慢}

课程明确指出，推理阶段的瓶颈往往是 memory-bound，而不是 compute-bound。生成新 token 时，模型需要反复访问之前所有 token 的 key/value；如果每次都从头计算，代价会很高。

\begin{figure}[H]
\centering
\includegraphics[width=0.78\textwidth]{pdf_pages/page-91.png}
\caption{KV caching 通过缓存历史 key/value，避免每生成一个 token 都重算全部过去。}
\end{figure}

KV cache 的基本思路就是把历史 token 的 key/value 存下来。这样生成第 \(t\) 个 token 时，只需计算最新 token 的投影，并与缓存交互。

\subsection{进一步压缩 KV 的存储成本}

缓存历史当然会节省重算，但也会消耗大量显存。于是课程把上一讲的 MQA/GQA 再拉回来，说明共享 key/value 头不仅是注意力近似技巧，也是 KV cache 压缩手段。

\begin{figure}[H]
\centering
\includegraphics[width=0.82\textwidth]{pdf_pages/page-104.png}
\caption{MQA/GQA 减少了需要缓存的 key/value 副本数，从而降低推理内存。}
\end{figure}

进一步地，PagedAttention 针对服务系统层的内存碎片问题，把 KV 以分页方式管理，减少浪费并提高吞吐。

\begin{figure}[H]
\centering
\includegraphics[width=0.88\textwidth]{pdf_pages/page-107.png}
\caption{PagedAttention 把 KV cache 做成分页管理，降低服务系统中的显存浪费。}
\end{figure}

课程还介绍了 latent attention 思路：不是缓存完整 K/V，而是缓存更低维的压缩表示，换取更小内存占用。

\begin{figure}[H]
\centering
\includegraphics[width=0.78\textwidth]{pdf_pages/page-110.png}
\caption{Latent attention 通过缓存压缩后的表示，进一步减小 KV cache 的内存占用。}
\end{figure}

\subsection{Speculative decoding 与 Multi-Token Prediction}

另一条加速路线不是压 KV，而是减少大模型逐 token 串行生成的等待时间。Speculative decoding 先让一个小 draft model 连续猜若干 token，再由大模型一次性验证。如果都被接受，就相当于一次前向推进多个 token。

\begin{figure}[H]
\centering
\includegraphics[width=0.82\textwidth]{pdf_pages/page-117.png}
\caption{Speculative decoding 用小模型起草 token，再由大模型并行验证。}
\end{figure}

课程特别强调它背后的概率论保证：通过接受-拒绝机制，可以让最终采样结果仍然匹配目标大模型分布，而不是仅仅得到一个启发式近似。

Multi-Token Prediction（MTP）则把“草稿模型”进一步吸收到同一模型内部，通过多个预测头一次预测多个未来 token，再由主头校验。它改变的是训练目标，因此和 speculative decoding 在原理上相近，但不完全等价。

\begin{figure}[H]
\centering
\includegraphics[width=0.82\textwidth]{pdf_pages/page-121.png}
\caption{MTP 在同一模型中加入多个预测头，以训练时的多 token 目标换取更快生成。}
\end{figure}

\subsection{本章小结}

推理优化的主线可以概括为两类：一类是少算或不重算，例如 KV cache、MQA/GQA、latent attention；另一类是更快推进生成步数，例如 speculative decoding 与 MTP。两类方法共同服务于同一个目标：在不明显损伤质量的前提下，提高吞吐、降低延迟。

\section{总结与延伸}

第三讲把“Transformer 课程”真正推进到了“LLM 课程”。回顾整讲，最值得沉淀的是五条主线：
\begin{enumerate}
\item \textbf{定义主线}：现代 LLM 通常是大规模 decoder-only 自回归模型。
\item \textbf{扩展主线}：MoE 通过稀疏激活扩大总参数规模，而不让单次激活成本同步爆炸。
\item \textbf{生成主线}：模型输出的是条件分布，真正的文本风格与稳定性很大程度上由解码策略决定。
\item \textbf{提示主线}：ICL、CoT、self-consistency 表明，不改参数也能显著改变任务表现。
\item \textbf{系统主线}：当模型足够大时，推理瓶颈进入内存与服务系统层，KV cache、PagedAttention、speculative decoding 等技巧因此成为核心基础设施。
\end{enumerate}

更重要的是，这些内容并不是分散的技巧库。它们一起刻画出今天 LLM 的真实形态：一方面，模型在结构上仍然站在 Transformer 的主干之上；另一方面，真正让它“可用”的，是大量围绕路由、采样、上下文和系统吞吐做出的工程决策。后续继续学习训练、对齐、agent 和评测时，这一讲给出的生成与推理框架会不断反复出现。

\end{document}
